%! Compressing Images by the SVD — Unit 7, Lesson 7.2 — Lesson Plan
\documentclass[10pt]{article}
\usepackage{linalg-article}
\usepackage{linalg-boxes}
\usepackage{graphicx}   % -article does NOT load graphicx; needed for thumbnails/screenshots
\graphicspath{{images/}}
\newcommand{\TallMath}[1]{$\displaystyle #1\rule[-1.4em]{0pt}{3.2em}$}
\newcommand{\uu}{\mathbf{u}}
\newcommand{\vv}{\mathbf{v}}
\newcommand{\T}{^{\top}}
\newcommand{\UnitNumberName}{Unit 7: The Singular Value Decomposition \quad}
\newcommand{\LessonNumberName}{Lesson 7.2: Compressing Images by the SVD}

\begin{document}
\begin{center}
    {\Large\bfseries \CourseName: \SchoolYear \\
    {\small\bfseries \UnitNumberName \LessonNumberName}}
\end{center}

\begin{tcolorbox}[colback=blush, colframe=burgundy, boxrule=0.9pt, arc=2mm,
    left=2mm, right=2mm, top=1.5mm, bottom=1.5mm]
    \textbf{Primary Objective:} Students can use the SVD to \emph{compress} a matrix. They build a rank-1
    matrix as an outer product $\uu\vv\T$ (a column times a row), write a matrix as a sum of ordered rank-1
    layers $A=\sigma_1\uu_1\vv_1\T+\sigma_2\uu_2\vv_2\T+\cdots$, and form the rank-$k$ approximation $A_k$ by
    keeping only the biggest layers. They can justify --- using the ``energy'' $\sigma^2$ and the storage
    count $1+m+n$ per layer --- why keeping the largest singular values gives the best low-storage copy of
    an image.
\end{tcolorbox}

\renewcommand{\arraystretch}{1.4}
\begin{skillbox}[Priority Ideas \& Skills]{goldbox}
\begin{tabularx}{\linewidth}{@{}X|X@{}}
\textbf{Priority Ideas \& Skills} & \textbf{Key Understandings (the ``why'')} \\[2pt]
\hline
\vspace{-6pt}
\begin{itemize}[leftmargin=1.1em, nosep, after=\vspace{2pt}]
  \item Build a rank-1 matrix as an outer product $\uu\vv\T$.
  \item Write $A$ as a sum of rank-1 layers $\sigma_i\uu_i\vv_i\T$, biggest $\sigma$ first.
  \item Form the rank-$k$ approximation $A_k$ by keeping the $k$ biggest layers.
  \item Measure the energy kept ($\sigma^2$) and the storage saved ($1+m+n$ per layer).
\end{itemize}
&
\vspace{-6pt}
\begin{itemize}[leftmargin=1.1em, nosep, after=\vspace{2pt}]
  \item A column times a row is the simplest nonzero matrix --- every column is a multiple of $\uu$.
  \item Ordering by $\sigma$ means the first layers carry the most of $A$.
  \item $A_k$ is the \emph{closest} rank-$k$ matrix to $A$ (Eckart--Young).
  \item Total energy $=\sigma_1^2+\sigma_2^2+\cdots$; big-$\sigma$ layers hold most of it, so a few store a near-perfect copy.
\end{itemize}
\end{tabularx}
\end{skillbox}

\begin{skillbox}[{Vocabulary, Concepts, \& Theorems}]{greenbox}
\renewcommand{\arraystretch}{1.3}
\begin{tabularx}{\linewidth}{@{}l X@{}}
\textbf{Outer product $\uu\vv\T$} & A column times a row; the result is a matrix whose every column is a multiple of $\uu$ (a \textbf{rank-1} matrix). \\
\textbf{Rank-1 layer $\sigma_i\uu_i\vv_i\T$} & One piece of the SVD --- a rank-1 outer product scaled by its singular value. \\
\textbf{SVD as layers} & $A=\sigma_1\uu_1\vv_1\T+\sigma_2\uu_2\vv_2\T+\cdots+\sigma_r\uu_r\vv_r\T$, ordered biggest $\sigma$ first. \\
\textbf{Rank-$k$ approximation $A_k$} & The sum of the first $k$ layers; the closest rank-$k$ matrix to $A$ (Eckart--Young). \\
\textbf{Energy} & The total size $\sigma_1^2+\sigma_2^2+\cdots$ (also the sum of squares of the entries); a layer's share is $\sigma_i^2$. \\
\end{tabularx}
\end{skillbox}

\begin{fixedskillbox}[Activate Prior Knowledge \& Spiral Review]{blush}
    The Warm-Up assembles the three pieces of today's idea on \emph{today's} matrix, so the Warm-Up answers
    become the opening of the notes: \textbf{(1)} a \emph{column times a row} (Unit~1 matrix multiply) ---
    the new outer-product object, with the observation that every column is a multiple of $\uu$;
    \textbf{(2)} the \emph{singular values} of the symmetric matrix $A=\begin{bsmallmatrix} 5 & 3\\ 3 & 5 \end{bsmallmatrix}$
    (Unit~6 / Lesson~7.1) --- because $A$ is symmetric and positive, $\sigma$'s equal the eigenvalues
    $8,2$ and $\uu_i=\vv_i$; and \textbf{(3)} a \emph{sum of squares}, showing $\sigma_1^2+\sigma_2^2=68$
    equals the sum of the squares of all four entries --- the ``energy'' that keeping big $\sigma$'s preserves.
\end{fixedskillbox}

\begin{skillbox}[Hook]{blush}
    Put up the number: a $1000\times1000$ grayscale photo is a matrix of \emph{one million} pixel values, and
    storing them all is expensive. Remind them of Lesson~7.1's promise: the SVD splits any matrix into
    \emph{ordered layers} $\sigma_1\uu_1\vv_1\T,\sigma_2\uu_2\vv_2\T,\dots$ Ask: \textbf{``If you could keep
    only a handful of layers and throw the rest away, which would you keep --- and how much of the picture
    would survive?''} That is exactly image compression: keep the biggest-$\sigma$ layers, drop the small
    tail, and store a near-perfect copy in a fraction of the space. Today we build the layers and count the
    savings.
\end{skillbox}

\begin{skillbox}[Lesson]{blush}
\begin{multicols}{2}
\textbf{The one new move: an outer product.} A column $\uu$ times a row $\vv\T$ is a whole matrix
$\uu\vv\T$, and every column of it is a multiple of $\uu$ --- the simplest nonzero matrix, called
\textbf{rank-1}. It stores in just $m+n$ numbers, not $mn$. \emph{Pose:} ``is $\uu\vv\T$ a number or a
matrix?'' (a matrix --- column times row, not row times column).

\textbf{The SVD is a stack of layers.} Multiplying out $A=U\Sigma V\T$ regroups $A$ into one rank-1 layer
per singular value: $A=\sigma_1\uu_1\vv_1\T+\sigma_2\uu_2\vv_2\T+\cdots$, biggest $\sigma$ first.

\columnbreak

\textbf{Worked spine $A=\begin{bsmallmatrix} 5 & 3\\ 3 & 5 \end{bsmallmatrix}$} ($\sigma_1=8,\sigma_2=2$;
$\uu_i=\vv_i$ since symmetric): layer~1 $=8\cdot\tfrac12\begin{bsmallmatrix} 1 & 1\\ 1 & 1 \end{bsmallmatrix}=\begin{bsmallmatrix} 4 & 4\\ 4 & 4 \end{bsmallmatrix}$,
layer~2 $=\begin{bsmallmatrix} 1 & -1\\ -1 & 1 \end{bsmallmatrix}$; they add back to $A$.

\textbf{Compress: keep the biggest.} $A_1=\begin{bsmallmatrix} 4 & 4\\ 4 & 4 \end{bsmallmatrix}$ is off by
only $\pm1$ and holds $\sigma_1^2/(\sigma_1^2+\sigma_2^2)=64/68\approx94\%$ of the energy. On a real image,
$20$ layers store $\approx4\%$ of the numbers and look almost identical.
\end{multicols}
\end{skillbox}

\begin{skillbox}[Explicit Instruction: Layer, Approximate, Count]{blush}
\begin{tabularx}{\linewidth}{@{}p{0.46\linewidth} X@{}}
\textbf{Steps.}
\begin{enumerate}[leftmargin=1.3em, nosep, after=\vspace{2pt}]
  \item Outer product: $\uu\vv\T$ (column times row) --- a rank-1 matrix.
  \item Layer: scale by $\sigma$ to get $\sigma\uu\vv\T$.
  \item Sum the layers to rebuild $A$ (biggest $\sigma$ first).
  \item Keep the first $k$ layers: $A_k$ (best rank-$k$ approximation).
  \item Count: energy kept $=\tfrac{\sigma_1^2+\cdots+\sigma_k^2}{\sigma_1^2+\cdots}$; storage $=k(1+m+n)$.
\end{enumerate}
&
\textbf{Worked} ($A=\begin{bsmallmatrix} 5 & 3\\ 3 & 5 \end{bsmallmatrix}$, $\sigma=8,2$).
$A_1=8\cdot\tfrac12\begin{bsmallmatrix} 1 & 1\\ 1 & 1 \end{bsmallmatrix}=\begin{bsmallmatrix} 4 & 4\\ 4 & 4 \end{bsmallmatrix}$;
dropped layer $\begin{bsmallmatrix} 1 & -1\\ -1 & 1 \end{bsmallmatrix}$.
Energy kept $=64/68\approx94\%$. For a $1000\times1000$ image, $20$ layers $=20(2001)\approx40{,}000$
numbers $\approx4\%$.
\end{tabularx}
\end{skillbox}

\begin{skillbox}[Active Monitoring]{redbox}
Circulate during the notes practice and Tier work. Watch for and correct:
\begin{itemize}[leftmargin=1.2em, nosep]
  \item computing an \emph{inner} product (row $\times$ column $=$ a number) instead of the \emph{outer} product (column $\times$ row $=$ a matrix);
  \item dropping the $\tfrac1{\sqrt2}\cdot\tfrac1{\sqrt2}=\tfrac12$ when scaling a layer (giving entries twice too big);
  \item forgetting to \emph{square} the singular values when computing the energy fraction;
  \item confusing ``rank-$k$'' (number of layers kept) with the size of the matrix.
\end{itemize}
Cold-call prompts: ``Is $\uu\vv\T$ a number or a matrix?'' ``Which layer carries the most --- how do you know?''
``What fraction of the energy does $A_1$ keep?'' ``How many numbers does one rank-1 layer cost?''
\end{skillbox}

\begin{skillbox}[Group Work \& Differentiation]{redbox}
\textbf{Compressing with rank-1 layers} activity (tiered; mirrors the notes).
\begin{multicols}{3}
\textbf{Tier R --- Remediate.} Build an outer product $\uu\vv\T$ and confirm it is rank~1; then scale one
into a layer $\sigma\uu\vv\T$ (warms up the Tier~A matrix).

\columnbreak
\textbf{Tier A --- Approaching.} Full compression of $B=\begin{bsmallmatrix} 4 & 2\\ 2 & 4 \end{bsmallmatrix}$:
build both layers, add them to recover $B$, form $B_1$, and judge it (keeps $36/40=90\%$ --- the best
rank-1 copy).

\columnbreak
\textbf{Tier E --- Extension.} The storage payoff: for a $600\times800$ photo, count full vs.\ rank-20
storage ($\approx5.8\%$) and \emph{justify} why the largest-$\sigma$ layers are the ones to keep (energy $\sigma^2$).
\end{multicols}
\end{skillbox}

\begin{skillbox}[Individual Work \& Assessment]{redbox}
\textbf{Exit Ticket} (independent, no notes): from a given $A=\begin{bsmallmatrix} 6 & 4\\ 4 & 6 \end{bsmallmatrix}$
with $\sigma=10,2$ and $\uu_i=\vv_i$, build the rank-1 approximation $A_1=\begin{bsmallmatrix} 5 & 5\\ 5 & 5 \end{bsmallmatrix}$;
compute the energy fraction it keeps ($100/104\approx96\%$) and judge whether it is a good copy; and a
\textbf{conceptual/justification check} --- why keep the largest singular values, and how does that save
storage? Collect and sort into ``got it / arithmetic slip / concept slip'' to decide whether 7.3 opens
with a re-cap.
\end{skillbox}

\begin{skillbox}[Reinforcement \& Extension]{goldbox}
\textbf{Homework:} outer products and their rank; a full two-layer decomposition of
$\begin{bsmallmatrix} 7 & 5\\ 5 & 7 \end{bsmallmatrix}$ with the rank-1 energy fraction ($144/148\approx97\%$);
a storage count on a $500\times500$ image ($\approx8\%$ for $20$ layers); justifications (keep biggest
$\sigma$; a layer costs $1+m+n$); and the $\sigma_2=0$ rank-1 case (one layer is exact).
\textbf{Extension:} a $\sigma$-spectrum ($10,6,3,1$) --- cumulative energy percentages and choosing the
smallest $k$ for $99\%$. \textbf{Preview:} Lesson~7.3, \emph{Principal Component Analysis} --- the same
``keep the biggest $\sigma$'' idea applied to data instead of pixels.
\end{skillbox}
\end{document}
